% Generated draft: requires mathematical and source review.
\documentclass[11pt]{article}
\usepackage[margin=25mm]{geometry}
\usepackage{fontspec}
\setmainfont{Latin Modern Roman}
\usepackage{amsmath,amssymb,mathtools,unicode-math}
\setmathfont{Latin Modern Math}
\usepackage{enumitem,needspace,xurl,hyperref}
\setlength{\parindent}{0pt}
\setlength{\parskip}{5pt}
\setlength{\emergencystretch}{3em}
\newcommand{\N}{\mathbb{N}}
\newcommand{\Z}{\mathbb{Z}}
\newcommand{\Q}{\mathbb{Q}}
\newcommand{\R}{\mathbb{R}}
\newcommand{\C}{\mathbb{C}}
\newcommand{\F}{\mathbb{F}}
\newcommand{\A}{\mathbb{A}}
\newcommand{\PP}{\mathbb{P}}
\newcommand{\E}{\mathbb{E}}
\newcommand{\eps}{\varepsilon}
\newcommand{\dd}{\,\mathrm{d}}
\begin{document}
\section*{UnsolvedMath ID: 20001100}
\textbf{Source problem code:} catalogue entry
\subsection{Definitions and mathematical statement}
\textbf{Recovered Problem 1.19 (M.-C. Chang).} The stored text has OCR-corrupted epsilon symbols and exponents. Following the aid's recovery, it asks whether, for every $\varepsilon>0$, there is $\tau=\varepsilon'>0$ such that for every sufficiently large integer $q$ and every $A\subset\Z/q\Z$ satisfying
\[
 |A+A|+|A\cdot A|<q^\varepsilon,
\]
either
\[
 |A|>q^{1-\tau},
\]
or some divisor $q_1\mid q$, $q_1>1$, has reduction image $|\pi_{q_1}(A)|\leq q^\tau$. The displayed small-image exponent uses the ambient $q$. In particular the stated divisor condition permits $q_1=q$, the identity projection. These choices are essential; a missing factor or a proper-divisor requirement would produce a different problem.

\textbf{Supplied resolution and a literal-normalization observation.} The aid attributes the assertion to Bourgain's arbitrary-modulus sum-product theorem and gives the explicit upper choice $\tau=\varepsilon$. It describes a relative sum-product lower bound under failure of the structural alternatives.

For the literal recovered normalization above, however, the same upper choice is immediate without that input: if $A\ne\varnothing$, translating $A$ by one of its elements shows $|A+A|\geq|A|$, hence $|A|<q^\varepsilon$. Taking $q_1=q$ supplies the second alternative. For empty $A$ that alternative also holds. Thus the reconstructed upper-existence statement needs no sufficiently-large threshold or Bourgain theorem. This does not dispute the named theorem; it flags that the exact OCR recovery and allowed projections must be checked before attaching a nontrivial parameter-transfer claim to this source.

\textbf{Reported exponent calibration.} For fixed $0<\varepsilon<1$, define $\tau_*(\varepsilon)$ as the infimum of exponents $\tau>0$ for which this recovered dichotomy holds for all sufficiently large moduli. The aid proves
\[
 \varepsilon/2\leq\tau_*(\varepsilon)\leq\varepsilon,
\]
and additionally that the dichotomy fails at $\tau=\varepsilon/2$ itself. The upper inequality agrees with the direct identity-projection observation. For prime moduli the only allowed divisor is $q_1=q$; the aid uses integer intervals to rule out every $\tau<\varepsilon/2$, and Ford's multiplication-table estimate to rule out the endpoint. Its exact interval parameters and Ford asymptotic are not supplied.

Nonattainment of the lower endpoint does not imply that the infimum is strictly larger than $\varepsilon/2$. The report also mentions a stronger trichotomy but gives no formula, so it cannot be stated precisely here. All calibration claims concern the displayed recovered absolute smallness condition, not an unprinted relative-small-doubling variant.
\subsection{Short English statement}
The OCR recovery gives an absolute small-sum/product condition and allows reduction modulo $q$ itself. With those exact choices, the reported upper exponent $\varepsilon'=\varepsilon$ follows directly from $|A+A|\geq|A|$. The aid separately reports a nontrivial lower-exponent obstruction, including failure at $\varepsilon/2$, whose details require source review.
\subsection{Sources}
\begin{itemize}
\item Local UnsolvedMath registry, permanent ID $20001100$, \texttt{catalogue entry}; canonical \texttt{aim-combinatorics-notes.json}, \texttt{notes[224]}.
\item AIM workshop ``Recent trends in additive combinatorics,'' PDF Problem 1.19, M.-C. Chang: \url{https://aimath.org/WWN/additivecomb/additivecomb.pdf}. Verification says the exact stored record is visibly OCR-corrupted; trailing ``Sumsets Problem Session'' starts another section.
\item Bourgain's arbitrary-modulus sum-product theorem and Ford's multiplication-table estimate, as named in the research summary; no complete references or precise theorem statements supplied.
\item Research attempt 1, \texttt{full\_solution}, run status \texttt{solved}, internal rigor \texttt{proved}, low novelty confidence; registry \texttt{SOLVED-IN-LITERATURE}. Explicitly machine-generated, not peer reviewed, and to be independently verified before citation. Original literature field empty.
\end{itemize}
\end{document}
